\chapter{Hasil dan Pembahasan}

\section{Pelaksanaan Kegiatan}
\lipsum[7]

\section{Pengujian Sistem}
Berikut adalah contoh pembuatan tabel profesional yang otomatis terdaftar di Daftar Tabel:

\begin{table}[htbp]
    \centering
    \caption{Hasil Pengujian Komponen}
    \label{tab:hasil_uji}
    \begin{tabular}{clc}
        \toprule
        \textbf{No} & \textbf{Skenario Uji} & \textbf{Status} \\
        \midrule
        1 & Rendering halaman utama & Berhasil \\
        2 & Pengambilan data API & Berhasil \\
        3 & Validasi form pengguna & Berhasil \\
        4 & Pengujian responsivitas UI & Berhasil \\
        \bottomrule
    \end{tabular}
\end{table}

Berdasarkan Tabel \ref{tab:hasil_uji}, seluruh fitur dasar berjalan tanpa hambatan.

\subsection{Implementasi Kode}
Potongan kode berikut menunjukkan implementasi dari komponen yang dibangun:

\begin{lstlisting}[language=JavaScript, caption={Komponen React Sederhana}, label={lst:react_code}]
import React, { useState } from 'react';

const UserProfile = () => {
    const [user, setUser] = useState('Dhiki');

    return (
        <div className="profile-container">
            <h1>Profil Pengguna</h1>
            <p>Selamat datang, {user}!</p>
        </div>
    );
};

export default UserProfile;
\end{lstlisting}

Listing \ref{lst:react_code} di atas merupakan contoh bagaimana state manajemen diterapkan.